% Generated draft: requires mathematical and source review.
\documentclass[11pt]{article}
\usepackage[margin=25mm]{geometry}
\usepackage{fontspec}
\setmainfont{Latin Modern Roman}
\usepackage{amsmath,amssymb,mathtools,unicode-math}
\setmathfont{Latin Modern Math}
\usepackage{enumitem,needspace,xurl,hyperref}
\setlength{\parindent}{0pt}
\setlength{\parskip}{5pt}
\setlength{\emergencystretch}{3em}
\newcommand{\N}{\mathbb{N}}
\newcommand{\Z}{\mathbb{Z}}
\newcommand{\Q}{\mathbb{Q}}
\newcommand{\R}{\mathbb{R}}
\newcommand{\C}{\mathbb{C}}
\newcommand{\F}{\mathbb{F}}
\newcommand{\A}{\mathbb{A}}
\newcommand{\PP}{\mathbb{P}}
\newcommand{\E}{\mathbb{E}}
\newcommand{\eps}{\varepsilon}
\newcommand{\dd}{\,\mathrm{d}}
\begin{document}
\section*{UnsolvedMath ID: 800564}
\textbf{Source problem code:} catalogue entry
\subsection{Definitions and mathematical statement}
Let \(G=(V,E)\) be a finite simple undirected graph with \(V=\{1,\ldots,n\}\). Its adjacency matrix is \(A\in\{0,1\}^{n\times n}\), where \(A_{ii}=0\) and, for distinct \(i,j\), \(A_{ij}=A_{ji}=1\) if and only if \(\{i,j\}\in E\). A clique is a subset \(C\subseteq V\) such that every two distinct vertices of \(C\) are adjacent. Define
\[
\omega(G)=\max\bigl\{|S|:S\subseteq V,\ \forall i,j\in S,\ i\ne j\Longrightarrow A_{ij}=1\bigr\}.
\]
Let \(\mathcal V\) denote the class of visibility graphs intended in the supplied source. The excerpt discusses visibility inside a polygon, but does not specify the geometric objects represented by vertices, the permitted polygon class, or the visibility convention. Thus \(\mathcal V\) remains an unresolved qualification of the statement.

\textbf{Computational problem (source Problem 38).} For every finite graph \(G\in\mathcal V\), given only its adjacency matrix \(A\), compute a clique \(C\subseteq V\) with \(|C|=\omega(G)\). The input is promised to be a visibility graph; no geometric realization is supplied. The excerpt specifies no running-time bound, so this is an optimization task whose intended algorithmic complexity target requires clarification.

The source states that a maximum clique corresponds to a largest empty convex polygon inside the corresponding polygon. The meaning of ``empty'' and the hypotheses needed for this correspondence are not supplied, so this observation is retained as a source assertion rather than used as an additional definition or an established equivalence. The source also reports that NP-hardness of this problem was not known in its 2010 discussion; this is a historical source qualification, not an independently verified current complexity assessment.
\subsection{Short English statement}
Given only the adjacency matrix of a graph promised to belong to the source's intended visibility-graph class, find a largest set of pairwise adjacent vertices. The exact visibility model and the desired running-time bound are unresolved in the supplied excerpt. This is a statement-only draft requiring human mathematical review.
\subsection{Sources}
\begin{itemize}
\item Local UnsolvedMath registry: permanent ID \(800564\), exact source code \texttt{catalogue entry}; title: \emph{Visibility Graphs of Points, Segments and Polygons --- Problem 38}. The supplied registry attributes the problem to Ghosh and Goswami, gives proposed year 2010, and labels its status ``open''. This status is supplied metadata, not independently verified.
\item Ghosh and Goswami, \emph{Unsolved Problems in Visibility Graphs of Points, Segments and Polygons} (2010), Problem 38, \url{https://arxiv.org/abs/1012.5187}. The supplied provenance records access on 2026-10-16 and extraction by \texttt{arxiv\_latex}; these are source-provided provenance details, not a claim of independent access or verification.
\item The excerpt cites reference labels \texttt{ar-clec-85}, \texttt{deo-secp-90}, and \texttt{es-mcmcpvg-00} for algorithms computing largest empty convex polygons. It says that the first two take a point set as input and the third takes a polygon as input. Full bibliographic entries are absent from the supplied record.
\item The excerpt attributes possible maximum-clique approaches using triangle-extendible ordering to Spinrad, reference label \texttt{s-grep-03}. Its description of that ordering as ``essentially a transitive orientation'' is a supplied informal characterization; neither a definition nor a full bibliographic entry is provided.
\end{itemize}
\end{document}
